\section*{結言}
この日の中心は、モデルの規模と実行可能性を巡る両方向の進展だった。Kimi K3.1 ProとMiniMax M3.1 Flashでは大規模MoEの提供形態とウェイト公開が注目され、Nemotron 3.5 Bolt E4Bでは小さなメモリ容量で使えるモデルの価値が議論された。オープン化の進展は、ライセンスと実装コストを含めた選択肢の比較を促している。

DeepInferを巡る最適化報告は、モデルを動かす推論スタックも継続的な競争領域であることを示す。一方、OpenAIの訓練環境隔離や脆弱性の話題は、モデルの能力が高まるほど周辺システムの安全性が重要になることを印象づけた。

なお、今回の7件には観測窓外の一次発表に基づく継続的な議論や、第三者による性能主張が含まれる。特に公開予定のウェイト、性能改善率、セキュリティ対策の効果は、今後の公式情報と再現性のある報告を待って評価する必要がある。
